% =============================================================================
% CAPÍTULO 1: INTRODUCCIÓN Y FUNDAMENTOS DEL CÓMPUTO
% =============================================================================

\chapter{Introducción y Fundamentos del Cómputo}
\label{ch:introduccion}

El estudio de los lenguajes de programación (PL) en el nivel de doctorado trasciende la sintaxis para enfocarse en los modelos matemáticos que permiten la computación. Como establece Sipser, la teoría de la computación busca responder qué es un computador mediante modelos idealizados como los \textbf{autómatas finitos} \cite{sipser2012introduction}. Este capítulo establece la base técnica del curso CIIC4030-066, vinculando la expresividad algorítmica con la eficiencia mecánica requerida en CISE.

\section{Motivación y Dominios de Aplicación}
\label{sec:1_1_motivacion_dominios}

\subsection{El lenguaje como herramienta de pensamiento algorítmico}

La razón primordial para estudiar PL es incrementar la capacidad del investigador para expresar ideas y algoritmos complejos. Para un físico, el lenguaje no es solo una librería; es el entorno donde se define la precisión del cálculo. Entender la mecánica de implementación permite elegir el lenguaje óptimo según el dominio, pasando de un enfoque de “herramienta” a uno de “diseño de sistemas” \cite{sebesta2019concepts}.

\begin{definicion}{Dominio Científico}
Dominio caracterizado históricamente por el manejo masivo de arreglos y cálculos de punto flotante de alta intensidad. Lenguajes como \textbf{Fortran} han dominado este espacio debido a su optimización para estas estructuras matemáticas \cite{sebesta2019concepts}.
\end{definicion}

\subsection{Especialización por Dominios de Aplicación}

Los lenguajes se clasifican según su optimización mecánica para tareas específicas \cite{sebesta2019concepts}:

\begin{itemize}
    \item \textbf{Inteligencia Artificial:} Enfocada en la manipulación de símbolos en lugar de números. El uso de lenguajes funcionales (LISP/Scheme) ha sido fundamental debido a su capacidad para manejar estructuras recursivas complejas.
    \item \textbf{Programación de Sistemas:} Donde la eficiencia es el criterio de diseño central. Lenguajes como \textbf{C} y \textbf{C++} permiten un control granular sobre el hardware, minimizando la latencia.
    \item \textbf{Aplicaciones Web:} Dominadas por lenguajes de marcado y scripts dinámicos para la interactividad en el cliente y servidor.
\end{itemize}

\begin{fisicaconexion}{}
En la simulación estocástica de reactores, la distinción entre dominios es crítica. Un algoritmo de transporte de neutrones que requiere \textbf{16 cifras significativas} de precisión no puede delegarse ciegamente a lenguajes con sistemas de tipado débil. La \textbf{confiabilidad (reliability)} del lenguaje garantiza que la propagación del error numérico sea predecible y controlable.
\end{fisicaconexion}

\subsection{Filosofía del Paradigma: Algoritmo vs. Gestión}

Una distinción técnica vital que surge en el curso es la carga cognitiva de la gestión de recursos. Mientras que en lenguajes como \textbf{Java o C}, el desarrollo se divide en un 50\% para el algoritmo y un 50\% para la gestión de memoria, los lenguajes de programación funcional (FPL) permiten que el investigador se enfoque en un \textbf{100\% en el algoritmo}.

\begin{ejemplo}{Criterio de Costo}
El costo de un lenguaje no solo es el tiempo de ejecución. Incluye el \textbf{costo de compilación}, el costo de entrenamiento del personal y, crucialmente, el costo de la \textbf{falta de confiabilidad}. En sistemas críticos como un reactor nuclear, un fallo en la gestión de memoria (por ejemplo, un \textit{memory leak}) representa un costo inaceptable comparado con la eficiencia bruta \cite{sebesta2019concepts}.
\end{ejemplo}

% =============================================================================
% SECCIÓN 1.2: CRITERIOS DE EVALUACIÓN DE LENGUAJES
% =============================================================================

\section{Criterios de Evaluación de Lenguajes}
\label{sec:1_2_criterios_evaluacion}

La selección de un lenguaje para el desarrollo de algoritmos científicos no es una decisión puramente estética; depende de una evaluación técnica que pondera la expresividad matemática frente al control mecánico de los recursos del sistema. Como señala Sipser, el estudio de la computación busca determinar con cuánta memoria y rapidez se puede resolver un problema \cite{sipser2012introduction}, lo que nos lleva a definir los siguientes criterios.

\subsection{Legibilidad y Capacidad de Escritura}

Para un físico computacional, la \textbf{legibilidad} (readability) es la facilidad con la que la lógica de un método numérico puede ser verificada por pares. La \textbf{capacidad de escritura} (writability) mide cuán natural es la traducción de una abstracción matemática a código \cite{sebesta2019concepts}.

\begin{itemize}
    \item \textbf{Simplicidad y Ortogonalidad:} Un lenguaje con alta ortogonalidad permite combinar un conjunto reducido de primitivas de formas predecibles. Esto es vital al construir tensores complejos en simulación física.
    \item \textbf{Tipos de Datos Adecuados:} La existencia de tipos nativos para números complejos o arreglos multidimensionales (como en \textbf{Fortran} o \textbf{Mathematica}) reduce la carga cognitiva del investigador \cite{sebesta2019concepts}.
\end{itemize}

\subsection{Confiabilidad (Reliability)}

La confiabilidad es el criterio más crítico en la simulación de reactores nucleares, donde un error de implementación puede tener consecuencias catastróficas \cite{sebesta2019concepts}.

\begin{itemize}
    \item \textbf{Verificación de Tipos (Type Checking):} El rastreo de errores de tipo en tiempo de compilación (\textit{static typing}) garantiza que las magnitudes físicas no se operen incorrectamente.
    \item \textbf{Manejo de Excepciones:} La capacidad de interceptar errores de punto flotante, como el \textit{overflow} en el cálculo de la reactividad $\eqk{\rho}$, permite una degradación elegante del sistema.
    \item \textbf{Aliasing:} La existencia de múltiples nombres para la misma celda de memoria. En algoritmos de transporte de neutrones, el \textit{aliasing} puede causar condiciones de carrera indeseadas en arquitecturas paralelas.
\end{itemize}

\begin{fisicaconexion}{Confiabilidad en la simulación de SDEs}
Al desarrollar métodos numéricos para \textbf{SDEs} (Ecuaciones Diferenciales Estocásticas), la confiabilidad del lenguaje asegura que la propagación del error numérico $\equ{\epsilon}$ no sea exacerbada por fallos en la gestión de memoria. Un lenguaje con \textit{Garbage Collection} (GC) determinista o inmutabilidad (como \textbf{Scala}) previene que el estado del sistema estocástico sufra mutaciones imprevistas durante la evaluación de trayectorias paralelas.
\end{fisicaconexion}

\subsection{Costo Computacional}

El costo de un lenguaje se analiza bajo la óptica de la complejidad asintótica de Sipser \cite{sipser2012introduction}. No solo incluye el tiempo de ejecución, sino también:

\begin{enumerate}
    \item \textbf{Costo de Compilación:} Crítico en fases de investigación donde el prototipado es frecuente.
    \item \textbf{Costo de Ejecución:} Determinado por la eficiencia del código objeto generado. Para un solver de PDEs, buscamos un tiempo de ejecución que se aproxime al límite teórico de la máquina.
    \item \textbf{Costo de Memoria:} El espacio $\eqk{S}(\eqk{n})$ requerido por el algoritmo.
\end{enumerate}

\begin{ejemplo}{Complejidad en la Práctica}
Un algoritmo de multiplicación de matrices masivas tiene un costo temporal de $O(\eqk{n}^3)$. Si el lenguaje elegido impone un sobrecosto (\textit{overhead}) por verificación de límites de arreglos en cada acceso, el factor constante en la notación \textit{Big-O} puede volver el algoritmo inviable para simulaciones de alta fidelidad, a pesar de ser teóricamente correcto \cite{sipser2012introduction}.
\end{ejemplo}


% =============================================================================
% SECCIÓN 1.3: MÉTODOS DE IMPLEMENTACIÓN
% =============================================================================

\section{Métodos de Implementación}
\label{sec:1_3_metodos_implementacion}

La transformación de un algoritmo desde su notación de alto nivel hasta su ejecución física puede realizarse mediante tres estrategias principales. Desde la perspectiva de Sipser, esta transformación es análoga a la simulación de una Máquina de Turing por otra, donde la eficiencia de la simulación depende del número de capas de abstracción \cite{sipser2012introduction}.

\subsection{Compilación Pura}

En el modelo de \textbf{compilación}, el programa fuente se traduce íntegramente a lenguaje de máquina antes de la ejecución. Este proceso maximiza el rendimiento, ya que permite optimizaciones globales intensivas durante la fase de \textit{back-end} \cite{sebesta2019concepts}.

\begin{definicion}{Compilador}
Un compilador es un traductor que mapea un programa $\eqk{P}$ escrito en un lenguaje fuente $\eqdom{L}_S$ a un programa equivalente $\equ{P}'$ en un lenguaje de destino $\eqdom{L}_D$ (usualmente código de máquina), tal que $\eqop{T}(\eqk{P}) \to \equ{P}'$ \cite{sipser2012introduction}.
\end{definicion}

\begin{fisicaconexion}{Compilación para alta precisión}
Para algoritmos de \textbf{cómputo de alta precisión} (16 cifras significativas), la compilación pura es preferible. Lenguajes como \textbf{Fortran} permiten al compilador realizar optimizaciones vectoriales y de flujo de datos que son críticas para minimizar la propagación del error numérico $\equ{\epsilon}$ en la resolución de \textbf{PDEs} masivas. La ausencia de un intérprete eliminando el \textit{overhead} asegura que el tiempo de ejecución se aproxime al límite de la CPU.
\end{fisicaconexion}

\subsection{Interpretación Pura}

La \textbf{interpretación} no genera un código objeto; en su lugar, un programa llamado intérprete actúa como un simulador de software para una máquina virtual cuyo “lenguaje de máquina” es el lenguaje de alto nivel \cite{sebesta2019concepts}.

\begin{itemize}
    \item \textbf{Ventajas:} Facilidad de depuración y gestión de errores en tiempo de ejecución.
    \item \textbf{Desventajas:} Ejecución significativamente más lenta (frecuentemente de 10 a 100 veces) debido a la decodificación repetitiva de instrucciones.
\end{itemize}

\begin{ejemplo}{La Máquina Universal de Turing}
Sipser define la \textbf{Máquina Universal de Turing} $\eqk{U}$ como el intérprete definitivo. $\eqk{U}$ toma como entrada la descripción de otra máquina $\eqk{M}$ y una cadena $\eqk{w}$, y simula el comportamiento de $\eqk{M}$ sobre $\eqk{w}$. Esta es la base teórica de todos los intérpretes modernos, desde \textbf{Python} hasta los motores de \textbf{Mathematica} \cite{sipser2012introduction}.
\end{ejemplo}

\subsection{Implementaciones Híbridas y JIT}

Muchos lenguajes modernos del curso, como \textbf{Scala} y \textbf{Go}, utilizan enfoques híbridos para equilibrar portabilidad y eficiencia \cite{sebesta2019concepts}.

\begin{acadtable}{tab:metodos_implementacion}{Métodos de implementación híbrida}
  \begin{tabularx}{\textwidth}{L{3cm} X}
    \toprule
    \headerrow
    \thead{Método} & \thead{Descripción Técnica} \\
    \midrule
    \textbf{Bytecode} & El código se traduce a un lenguaje intermedio (como en la JVM para Scala) que luego es interpretado o compilado bajo demanda \cite{sebesta2019concepts}. \\
    \addlinespace
    \textbf{JIT (Just-In-Time)} & Compilación selectiva de fragmentos de código intermedio a código de máquina durante la ejecución para mejorar el rendimiento de bucles calientes. \\
    \addlinespace
    \textbf{AOT (Ahead-Of-Time)} & Compilación previa al tiempo de ejecución, común en \textbf{Go}, que genera binarios estáticos de alta velocidad \cite{sebesta2019concepts}. \\
    \bottomrule
  \end{tabularx}
\end{acadtable}

\begin{fisicaconexion}{Híbridos en IA Generativa}
En la investigación de \textbf{IA Generativa}, el modelo híbrido es común. Mientras que la lógica de alto nivel (el grafo de la red neuronal) se define en un entorno interpretado como Python para facilitar la experimentación, los \textit{kernels} de cómputo estocástico y las operaciones de tensores $\eqten{T}$ se ejecutan en código compilado (C++/CUDA) para mantener la eficiencia necesaria en el entrenamiento de modelos de difusión.
\end{fisicaconexion}


% =============================================================================
% SECCIÓN 1.4: LA JERARQUÍA DE LA COMPUTACIÓN Y EL PIPELINE
% =============================================================================

\section{La Jerarquía de la Computación en el Pipeline de Traducción}
\label{sec:1_4_jerarquia_computacion}

El diseño de un compilador moderno se fundamenta en la correspondencia biyectiva entre las clases de lenguajes formales y las máquinas abstractas que los reconocen. Esta clasificación, conocida como la \textbf{Jerarquía de Chomsky}, dicta la complejidad $\eqop{O}(\eqk{f}(\eqk{n}))$ de cada fase del proceso de traducción \cite{sipser2012introduction}.

\subsection{Lenguajes Regulares y el Lexer (Tipo 3)}

La fase inicial de compilación, el \textit{Lexer} o escáner, se encarga de agrupar caracteres en unidades con significado léxico llamadas \textit{tokens}. Matemáticamente, este proceso se modela mediante \textbf{Autómatas Finitos Deterministas (DFA)}, los cuales tienen una memoria extremadamente limitada \cite{aho2006compilers,sipser2012introduction}.

\begin{definicion}{Lenguaje Regular}
Un lenguaje es regular si existe un autómata finito que lo reconoce. En el contexto de un compilador, las expresiones regulares definen la estructura de los identificadores, constantes numéricas y palabras clave \cite{aho2006compilers,sipser2012introduction}.
\end{definicion}

\begin{observacion}{Análisis léxico y procesamiento carácter a carácter}
Como señala el texto clásico de compiladores \cite{aho2006compilers}, la eficiencia del análisis léxico es vital porque es la única fase del compilador que debe examinar cada carácter del código fuente uno por uno, y cualquier ineficiencia se multiplica por el tamaño total del programa.
\end{observacion}

\subsection{Lenguajes Libres de Contexto y el Parser (Tipo 2)}

Una vez generados los tokens, el \textit{Parser} verifica la estructura jerárquica del programa (por ejemplo, que cada paréntesis abierto tenga uno de cierre). Esta fase requiere una memoria de tipo LIFO (\textit{Last-In, First-Out}), formalizada por el \textbf{Autómatas de Pila (PDA)} \cite{aho2006compilers,sipser2012introduction}.

\begin{itemize}
    \item \textbf{BNF (Backus-Naur Form):} Es la notación estándar para describir gramáticas libres de contexto. Fue introducida por John Backus (creador de \textbf{Fortran}) y Peter Naur para la especificación de ALGOL 60.
    \item \textbf{Derivación:} El proceso de aplicar reglas gramaticales hasta obtener una sentencia compuesta solo por terminales $\eqk{T}$ \cite{aho2006compilers}.
\end{itemize}

\begin{fisicaconexion}{Consistencia dimensional y análisis sintáctico}
Para un físico, el análisis sintáctico es análogo a la verificación de la \textbf{consistencia dimensional} en una ecuación. Si la gramática del lenguaje (especificada en BNF) no permite la operación entre ciertos tipos de datos, el parser detectará el error estructural antes de que el algoritmo intente computar resultados físicamente incoherentes.
\end{fisicaconexion}

\subsection{Computabilidad Universal y Ejecución (Tipo 0)}

Finalmente, la ejecución del código objeto resultante ocurre en un entorno equivalente a una \textbf{Máquina de Turing}. Aquí es donde el paradigma funcional, inspirado en el \textbf{Cálculo Lambda} de Alonzo Church, ofrece una alternativa a la computación basada en estados de Turing \cite{sipser2012introduction}.

\begin{acadtable}{tab:jerarquia_chomsky}{Jerarquía de Chomsky en el pipeline de compilación}
  \begin{tabularx}{\textwidth}{C{2.5cm} L{3.5cm} X}
    \toprule
    \headerrow
    \thead{Tipo} & \thead{Lenguaje} & \thead{Máquina Abstracta} \\
    \midrule
    \textbf{Tipo 3} & Regular & Autómata Finito (Lexer) \\
    \addlinespace
    \textbf{Tipo 2} & Libre de Contexto & Autómata de Pila (Parser) \\
    \addlinespace
    \textbf{Tipo 1} & Sensible al Contexto & Autómata Linealmente Acotado \\
    \addlinespace
    \textbf{Tipo 0} & Recursivamente Enumerable & Máquina de Turing / Cálculo Lambda \\
    \bottomrule
  \end{tabularx}
\end{acadtable}

\subsection{Limitaciones de la Abstracción}

Entender esta jerarquía permite al investigador en \textbf{CISE} discernir por qué ciertos problemas son intratables. Como advierte Michael Sipser, la teoría de la computación nos proporciona las herramientas para no intentar resolver mediante algoritmos aquello que es intrínsecamente \textbf{indecidible} \cite{sipser2012introduction}.

\begin{ejemplo}{Complejidad del Análisis Sintáctico}
El análisis sintáctico de un lenguaje de programación debe realizarse en tiempo lineal o cuasi-lineal $\eqop{O}(\eqk{n})$. Si se utilizara un algoritmo de propósito general para cualquier gramática libre de contexto, la complejidad ascendería a $\eqop{O}(\eqk{n}^3)$, lo cual sería inaceptable para procesar grandes bases de datos de simulaciones estocásticas \cite{aho2006compilers}.
\end{ejemplo}

\printbibliography[heading=subbibliography]
